\VerdierPhysicalPageAnchor{0073}%
\VerdierSemanticAnchor{ch1-2-3-4}
\noindent\textbf{2.3.4.}\quad
Prenons pour application $\varphi$ l'unique application
$[n]\xrightarrow{\varphi}[1]$ et faisons opérer le groupe $\mathbb Z[n]$
sur $\operatorname{comp}(C)$ par l'intermédiaire de
$\varphi\colon\mathbb Z[n]\longrightarrow\mathbb Z$. Il existe alors,
d'après (1.5.2), un seul $\mathbb Z[n]$-foncteur tordu par le cocycle de
Koszul :

\VerdierSemanticAnchor{ch1-2-3-4-eq1}
\begin{equation}
  \left(\int_n^\Sigma,p_n\right)\colon
  \operatorname{comp}^{n}(C)\longrightarrow\operatorname{comp}(C)
  \tag{2.3.4.1}
\end{equation}

\VerdierSemanticAnchor{ch1-2-3-4-characterization}
tel que :
\begin{equation*}
  p_n(e_i)=p_{i,\varphi}.
\end{equation*}

\VerdierSemanticAnchor{ch1-2-3-4-notation}
Dorénavant, la notation $\int_n^\Sigma$ désignera le foncteur
$\int_n^\Sigma$ muni de la structure de $\mathbb Z[n]$-foncteur tordu
décrite ci-dessus.

\VerdierSemanticAnchor{ch1-2-3-5}
\noindent\textbf{2.3.5.}\quad
Les considérations développées aux alinéas (2.3.2), (2.3.3) et (2.3.4)
s'appliquent, modulo les modifications évidentes, aux foncteurs
$\int_\varphi^\Pi$.

\VerdierSemanticAnchor{ch1-2-3-5-summability-restriction}
Lorsque les sommes et produits dénombrables ne sont pas nécessairement
représentables dans $C$, les résultats ci-dessus s'appliquent aux foncteurs
$\int_\varphi$ en se restreignant à des complexes possédant des propriétés
de sommabilité convenables.

\VerdierSemanticAnchor{ch1-2-4}
\subsection{Extension des foncteurs aux complexes}

\VerdierSemanticAnchor{ch1-2-4-1}
\noindent\textbf{2.4.1.}\quad
Il existe un foncteur canonique :

\VerdierSemanticAnchor{ch1-2-4-1-eq1}
\begin{equation}
  \iota_n\colon
  \operatorname{comp}^{n}(C)^\circ
  \longrightarrow
  \operatorname{comp}^{n}(C^\circ)
  \tag{2.4.1.1}
\end{equation}
défini par les formules suivantes :

\VerdierSemanticAnchor{ch1-2-4-1-object-formula}
\begin{equation*}
  (\iota_nX)^\sigma=X^{-\sigma},
  \qquad X\in\operatorname{Ob}\bigl(\operatorname{comp}^{n}(C)\bigr),
\end{equation*}

\VerdierSemanticAnchor{ch1-2-4-1-eq2}
\begin{equation}
  d_{\iota_nX}^{\sigma,i}=d_X^{-\sigma-e_i,i},
  \tag{2.4.1.2}
\end{equation}

\VerdierSemanticAnchor{ch1-2-4-1-morphism-formula}
\begin{equation*}
  (\iota_nf)^\sigma=f^{-\sigma},
  \qquad f\in\operatorname{Fl}\bigl(\operatorname{comp}^{n}(C)\bigr).
\end{equation*}

\VerdierSemanticAnchor{ch1-2-4-1-graded-isomorphism}
Le foncteur $\iota_n$ est un isomorphisme de $\mathbb Z[n]$-catégories
(la structure de $\mathbb Z[n]$-catégorie de
$\operatorname{comp}^{n}(C^\circ)$ est définie par (2.3.1) ; la structure
graduée de $\operatorname{comp}^{n}(C)^\circ$ est définie par (2.3.1) et
(1.3.9)).

\VerdierSemanticAnchor{ch1-2-4-1-sigma-opposite-setup}
Supposons que les sommes directes dénombrables soient représentables dans
$C$ et soit $\varphi\colon[n]\longrightarrow[m]$ une application.
Désignons par :

\VerdierSemanticAnchor{ch1-2-4-1-eq3}
\begin{equation}
  \left(\int_{\varphi,C}^{\Sigma}\right)^\circ\colon
  \operatorname{comp}^{n}(C)^\circ
  \longrightarrow
  \operatorname{comp}^{m}(C)^\circ,
  \tag{2.4.1.3}
\end{equation}
